We present the first machine-checked correctness proof of the
OpenZeppelin reentrancy-guard pattern against a Lean 4 state-machine
model of production-deployed Solidity source with a composition
meta-theorem spanning multiple production protocols. All thirteen
theorems are machine-checked in Lean 4 with \textbf{zero \texttt{sorry},
zero user-introduced axioms, and an axiom footprint bounded by
\texttt{{[}propext{]}} across the entire corpus} --- propositional
extensionality, a standard classical axiom in Lean 4's mathlib4;
per-function inner lemmas are kernel-only, master/wrapper theorems carry
\texttt{{[}propext{]}}-only records, and the Layer 6-D capstone composes
exactly the union of the prior-layer records by direct conjunction. The
corpus is gated under continuous integration across four parallel
verification blocks that re-check each theorem and its axiom record on
every push.

Smart contract reentrancy has caused over US\$500M in cumulative
documented losses since 2016, with the DAO 2016 attack alone draining
≈3.6M ETH and forcing a contentious hard fork that split Ethereum into
two persistent chains. The OpenZeppelin \texttt{ReentrancyGuard} mutex
pattern has emerged as the de facto defense across most production DeFi
protocols, yet no prior formal verification effort has established its
\emph{discriminating power} --- the property that the guard pattern
blocks attacks against vulnerable instances, preserves correct execution
for non-attacking transactions, and distinguishes structurally-adjacent
safe and vulnerable variants. Prior work has formalized either guard
correctness on toy contracts or attack feasibility on isolated
vulnerable instances, but not both directions plus boundary cases
against production-deployed source.

The verification covers three production protocol instantiations --- DAO
2016, Compound v2 cToken family, and Aave V3 \texttt{flashLoan} ---
together with one constructed minimal-diff mutant of Aave V3's
production \texttt{flashLoan} (\texttt{flashLoanVulnerable}) authored to
isolate a single security-critical structural difference for
discriminating-power isolation. We apply mutation-testing methodology to
formal verification by constructing this minimal-diff mutant, enabling a
controlled experiment that naturally-occurring near-misses rarely
provide. The resulting tridirectional structure pairs (a) a
negative-instance attack reproduction of the DAO 2016 vulnerable
pattern, (b) a positive-instance correctness proof against the Compound
v2 cToken family, and (c) a boundary-case proof distinguishing Aave V3's
safe-by-design \texttt{flashLoan} (CEI-correct) from the
\texttt{flashLoanVulnerable} mutant. A single capstone meta-theorem
composes the three protocol instantiations under a \emph{no-retrofit
composition discipline} (each protocol-instantiation proof was sealed
prior to capstone authoring, with no underlying-proof modifications
during composition), establishing a machine-checked composition
meta-theorem demonstrated at the first cross-protocol stress test
(Compound v2 → Aave V3, within the lending-family domain; broader-family
portability is future work, §9.4).

We release the full Lean 4 source, CI gating configuration, and
reproduction commands at
\texttt{https://github.com/rayiskander2406/qanary-contracts} with
reproducibility verified at tagged commit \texttt{v1.6-phase7-closure}
(the post-audit content seal; the substantive proof substrate is
independently reproducible at \texttt{v1.3-layer6-closure} per Appendix
A.3).
